\subsection{A decision-preservation test for learned abstraction}
\label{sec:C05:decisiontest}
An abstraction can be evaluated more directly than through a generic probe score. Let $X$ be the complete permitted observation and let $Y$ indicate a binary action requirement, with $Y$ balanced. Construct two environments whose broad appearance is the same but whose required action differs because of a localized signal. If the encoder maps both observations to exactly the same representation and no additional evidence is available, the best downstream classifier has accuracy at most one half. This follows from the state-aliasing argument in Chapter~\ref{ch:C03}. With approximate rather than exact collisions, the relevant object is the separability of the induced representation distributions under the actual decoder class.

For unrestricted binary testing with equal priors, the optimal success probability is $(1+\operatorname{TV}(P_0,P_1))/2$, where $P_y$ is the representation distribution conditional on $Y=y$. The identity follows by choosing the decision region that maximizes $P_1(A)-P_0(A)$. This is a standard probability deduction, used here to specify a diagnostic rather than to claim a new representation theorem. It explains why a low average reconstruction error or a high semantic retrieval score cannot alone certify preservation of a rare action-critical signal: those metrics need not control the total variation of these conditional distributions.

In practice, total variation in high-dimensional learned representations is not directly known. A fixed classifier with population accuracy $A_d$ gives $\operatorname{TV}(P_0,P_1)\geq\max\{0,2A_d-1\}$, not an upper bound on the information present. A finite held-out accuracy estimate needs a one-sided lower confidence bound $A_L$, valid for the sampling design, before reporting $\max\{0,2A_L-1\}$ as a statistical lower bound. Tuning on those test outcomes invalidates that independent-test interpretation. A failed probe can therefore motivate stronger probes or a direct action test, while a successful probe still leaves policy access and calibration unresolved. The experiment should report encoder stochasticity, preprocessing, sample size, decoder capacity, and the precise intervention defining $Y$. Those details determine what conclusion can be drawn from the observed accuracy.

The resulting design principle is to preserve a recoverable path from task-critical observations to execution, while allowing abstraction of genuine nuisances. Multiple channels can provide complementary evidence, but agreement between two models trained on similar data is not independence. A robust agent records which channel supplied a fact, checks freshness before action, and seeks another observation when the distinction remains unresolved. Representation learning contributes powerful statistical structure; the surrounding system must still connect that structure to the exact state and action semantics of the task.

\subsection{Directed progress and finite-budget reachability}
\label{sec:C05:plangeometry}
RC-aux changes two interfaces that one-step JEPA prediction leaves unspecified: recursive prediction over the planning horizon and a directed reachability estimate conditioned on remaining steps. Its Section 3 constructs positive and negative labels from observed trajectory offsets and also supplies predicted-latent inputs to the reachability head. Temporal-Distance JEPA's current second version instead calibrates a nonnegative asymmetric cost from demonstration step counts, with a metric-residual head and horizon-matched rollout supervision~\cite{rewrite_rcaux2026,rewrite_tdjepa2026}. A distance between two embeddings is therefore no longer assumed to be the planning objective merely because the predictor was trained in that space.

The supervision has a precise epistemic limit. In deterministic state dynamics, let $D^*(s,g)$ be the minimum number of feasible actions from $s$ to $g$, with value $+\infty$ when none exists. An observed valid segment from $s_i$ to $s_j$ of length $j-i$ proves $D^*(s_i,s_j)\leq j-i$. It does not prove equality. Consequently, labeling the pair unreachable at budget $h<j-i$ introduces possible false negatives whenever an unobserved shortcut exists. Cross-trajectory negatives are weaker still: two independently collected trajectories may pass through mutually reachable states. These are deductions from the labeling construction, rather than objections that invalidate every useful learned ranking.

Directionality is independently necessary when actions are irreversible. If a transition deletes an object and no action restores it, $D^*(s,g)$ can be finite while $D^*(g,s)$ is infinite. Every symmetric embedding distance assigns the same value in both directions. An asymmetric head removes that representational restriction, but does not establish calibrated reachability on unseen states. To distinguish the interventions, hold the encoder and rollout operator fixed while comparing symmetric and directed heads, then hold the scoring head fixed while changing recursive prediction training. Include observed loops and deliberately collected shortcut trajectories to measure label ambiguity.

For a stateful agent, the budget must count executable operations rather than an unspecified number of latent updates. A plan with three service calls and a required approval can be unreachable within three permitted actions even if its visual endpoint is nearby. Test changes in permission and resource availability while preserving the displayed scene. A reachability head that sees only pixels cannot resolve a hidden authorization change without another state channel. The strongest diagnostic therefore reports action feasibility, goal completion, and predicted reachability separately, including false acceptance of infeasible plans. Neither source's trajectory-based supervision supplies a general certificate for that broader execution contract.

\subsection{Latent projection and the difference between predictability and control}
\label{sec:C05:projection}
Latent Projection for Latent Planning fits an action-free linear dynamics operator to frozen representations and changes the goal cost through a selected subspace. Its Sections 3--4 assume joint Gaussianity, stationary standard-normal marginals, Markov dynamics, and a normal transition matrix~\cite{rewrite_latentprojection2026}. Writing $z'=Az+\eta$ with independent zero-mean innovation, stationarity gives
\begin{equation}
 \operatorname{Cov}(\eta)=I-AA^\top\succeq0.
\end{equation}
This follows by taking the covariance of the transition, not from an empirical assurance that a regularizer has made real embeddings exactly Gaussian. Normality further yields $A^h(A^h)^\top=(AA^\top)^h$, so one orthonormal basis diagonalizes the innovation covariance at every positive integer horizon. Without normality, the preferred subspace can change with horizon.

The pinned implementation makes several consequential choices visible. \texttt{NormalDrift} in \texttt{module.py} parameterizes contraction and rotation blocks; its evaluation mode zeros configured static line blocks. \texttt{LP2GoalMSE} in \texttt{cost.py} selects left singular vectors of the evaluation operator and compares terminal features after projection~\cite{rewrite_lp2_repo}. Thus the exact choice of dynamic planes, static lines, retained rank, and training versus evaluation mode belongs in the experimental configuration. A description that mentions only top singular vectors omits an implemented structural filter.

A counterexample limits any universal identification of predictable features with controllable goal progress. Let $z=(x,y)\in\mathbb R^4$, where $x,y\in\mathbb R^2$, and define hypothetical dynamics
\begin{align}
 x'&=rR x+\sqrt{1-r^2}\,\xi,\\
 y'&=a,
\end{align}
with $0<r<1$, orthogonal two-dimensional rotation $R$, and independent $\xi\sim\mathcal N(0,I_2)$. Under behavior actions $a\sim\mathcal N(0,I_2)$ independent of the current state, stationary marginals are standard Gaussian and the action-free fitted operator is $A=\operatorname{diag}(rR,0_2)$. It is normal and its most predictable two-dimensional subspace is $x$. Yet actions directly set $y$, so a goal depending on $y$ is reachable in one step while a projected cost retaining only $x$ cannot distinguish candidate actions by that goal.

This constructed example is not a reported failure of the paper's evaluated models. It proves that the stated statistical assumptions alone do not guarantee decision sufficiency, even when the retained predictable features are dynamic planes rather than static lines. A falsifiable test introduces an autonomous visual distractor and a less behavior-predictable but controllable target, then compares full latent cost, fitted projection, random projection, and an action-conditioned relevance control at matched search budget. Vary the behavior policy used to fit the projection. A change in the selected subspace can reveal dependence on exploration statistics rather than an invariant property of environment controllability.
